% Rubric: Experimental Design.
% Specify scenarios, parameter ranges, outcome measures and, where relevant,
% repeated runs and comparisons, with enough detail for the experiments to be
% repeated.

\section{Experimental design}
\label{sec:experiments}

Five runner families produce the evidence, each from one command in the
repository root. They load the frozen graph, apply the model of
Section~\ref{sec:model}, and write their tables and figures under
\texttt{results/} and \texttt{figures/}. Table~\ref{tab:design} lists each
family's scenarios, runner defaults and outcome measures.

\begin{table}[t]
  \centering
  \scriptsize
  \setlength{\tabcolsep}{4pt}
  \caption{Scenarios, runner defaults and outcome measures.}
  \label{tab:design}
  \begin{tabular}{@{}p{0.15\linewidth} p{0.49\linewidth} p{0.33\linewidth}@{}}
    \toprule
    Command & Scenarios and runner defaults & Outcome measures \\
    \midrule
    RQ1: \texttt{make percolation} &
    Random and targeted node and edge removal ranked by degree, betweenness
    and demand flow; fractions 0.00 to 0.90 in 0.01 steps plus 1.00, 100 seeds
    each; every station once as trigger over 81 tolerances ($\alpha = 0$ to
    $2$, step 0.025), static capacity rule &
    Seed-mean GCC fraction and isolates; susceptibility peak; GCC $= 0.5$
    crossing; containment tolerance; predictor strength and rank stability \\
    \addlinespace
    RQ2: \texttt{make cascade} &
    Static sweep $\alpha = 0$ to $2$ in 0.025 steps with \texttt{equal} and
    \texttt{capacity} rules; maximum-load, maximum-degree and 300 seeded
    random triggers. A reduced 10-point grid adds the frequency and demand
    load modes and dynamic rerouting. 2{,}000 random triggers at
    $\alpha = 0.1, 0.15, 0.2, 0.3$; line closures as one batch &
    Failed fraction; avalanche size, duration and span; containment
    $\alpha^*$; paired rule, load-mode and rerouting differences \\
    \addlinespace
    RQ3: \texttt{make recovery} &
    21 closures: 8 exclusive-station lines, 8 interchanges, 5 seeded random
    sets of 5. Five strategies with a 2-link budget over 4 manual and 266
    GTFS-derived candidate pairs; layered cascade $\alpha = 0$ to $2$ in 0.05
    steps; focus checks at 0 and 0.2 &
    Served OD fraction and endpoint-sum demand over 3{,}655 terminal pairs;
    secondary rail closures and recovery rounds; travel-time penalty
    conditional on reachable pairs \\
    \addlinespace
    RQ4: \texttt{make demand} &
    Topology (betweenness loads, $K = L0$) against demand
    (\texttt{am\_peak\_stops} loads, frequency-scaled $K$); 81 tolerances
    from 0 to 2 in 0.025 steps; maximum-load trigger; capacity rule &
    Load and rank changes; Spearman correlation; top-5, top-10 and top-20
    overlap; containment-tolerance shift \\
    \addlinespace
    RQ4: \texttt{make uncertainty} &
    1{,}000 random draws per removal fraction (0 to 1, step 0.05); 95\%
    percentile bootstrap with 2{,}000 resamples; nested seed prefixes 25 to
    2{,}000; alpha grids from 0.5 to 0.025; power-law test with 500 refitted
    simulations at $\alpha = 0.1, 0.15, 0.2, 0.3$ &
    Pointwise interval half widths; seed convergence; grid-induced
    containment shift; tail eligibility and rejection probability \\
    \bottomrule
  \end{tabular}
\end{table}

All runs use master seed 0 and derived child seeds, so paired conditions share
random draws while repeated draws remain separate observations. Each runner
writes one row per replicate through \texttt{experiments.save\_table}, with a
\texttt{.meta.json} sidecar holding the seed, parameters, package versions and
SHA-256 hashes of inputs and code, plus a family manifest that hashes the
artifacts; \texttt{--quick} variants are smoke tests in separate directories.

The pipeline is runner, stored rows, statistics, report. The runners call the
packaged engines and store raw per-seed rows under \texttt{results/<family>/};
the P1.2 helpers summarise them as means and percentile intervals;
\texttt{make figures} checks the tables against their sidecars and manifests,
then writes the report figures and the macros in
\texttt{sections/00-results-values.tex}; and Section~\ref{sec:results} cites
those frozen files. Notebooks 01 to 04 (issues P3.1 to P3.4) re-execute the
same calls interactively. \texttt{make check} runs the test suite and
\texttt{make reproduce} chains \texttt{data}, \texttt{check},
\texttt{figures} and \texttt{notebooks}; each family regenerates with its
command in Table~\ref{tab:design}. \path{run_uncertainty.py} verifies the
committed P2.2 manifest before running, so P2.2 is regenerated first if its
inputs change.